\documentclass[10pt,a4paper]{amsart}
\usepackage[margin=19mm]{geometry}
\usepackage{amsmath,amssymb,mathtools}
\usepackage[scr=rsfs,cal=euler]{mathalfa}
\usepackage{xfrac}
\usepackage[unicode,hidelinks,bookmarks=false]{hyperref}
\newcommand{\ie}{i.e.\ }
\newcommand{\cf}{cf.\ }
\newcommand{\resp}{respectively\ }
\newcommand{\Subsection}{\subsection}
\providecommand{\nobreakdash}{\nobreak\hbox{-}}
\setlength{\emergencystretch}{2em}
\multlinegap0pt
\usepackage{amssymb}
\usepackage[all]{xy}
\usepackage{mathtools}
\usepackage{xcolor}
\newtheorem{lemma}{Lemma}
\newcommand{\bA}{\mathbb{A}}
\newcommand{\cA}{\mathcal{A}}
\newcommand{\cS}{\mathcal{S}}
\newcommand{\rd}{\mathrm{d}}
\title{The global Gan--Gross--Prasad conjecture for Fourier--Jacobi periods on unitary groups II: Comparison of the relative trace formulae}
\author{Paul Boisseau, Weixiao Lu, Hang Xue}
\date{Source: arXiv:2601.01727v1 (CC BY 4.0)}
\begin{document}
\maketitle

\noindent\emph{Source and licence.} The global Gan--Gross--Prasad conjecture for Fourier--Jacobi periods on unitary groups II: Comparison of the relative trace formulae. Version arXiv:2601.01727v1 (CC BY 4.0), distributed by arXiv under the Creative Commons Attribution 4.0 International licence, http://creativecommons.org/licenses/by/4.0/, as recorded in the arXiv licence metadata for this exact version and shown on the abstract page https://arxiv.org/abs/2601.01727v1. Full licence notice: \texttt{licenses/arxiv-2601.01727-CC-BY-4.0.txt}.\par\smallskip

\noindent\textit{Editorial scope.} Counted unit: the article's lemma lem:geometric\_linear\_convergence (source0.tex lines 921--939). The article's complete original proof of it is delivered with it (source0.tex lines 942--1033, one \texttt{proof} environment of 92 lines). No mathematics is written, repaired or re-derived: the statement and the proof are the article's own lines, in the article's own notation, and no retained line is substituted or re-flowed.  No further source-local context is needed: every cross-reference of every retained line resolves inside the two delivered spans.  Nothing was omitted from any retained span, so the package carries no omission record.  No retained line contains a citation, so the wrapper carries no bibliography and no bibliography entry.  No retained line contains a figure environment, an image-inclusion command or drawing code, so no image is delivered and the package declares no asset dependency.  Editorial formatting only, in addition: an AMS \texttt{amsart} preamble that restates the article's own declarations and macros used by the retained lines, namely the environments \texttt{lemma} on separate counters, together with the standard packages those lines need.  The retained environments print in this document as Lemma 1 in order of appearance, whereas the article prints the same items with the numbering of its own full document.  The complete native source of the article as distributed in the arXiv e-print for this exact version is bundled unchanged as \texttt{sources/192101/source0.tex}.
\begin{lemma}   \label{lem:geometric_linear_convergence}
For all $\alpha \in \cA(F)$ and $f^+ \in \cS(G^+(\bA))$, the expression
    \[
    \int_{[H]} \int_{[G_n']}
    \sum_{ \substack{ \gamma \in P_H(F) \backslash H(F) \\
    \delta_2 \in P_2'(F) \backslash G_n'(F) }}
    \sum_{m \in M_{\alpha}^+(F)} \int_{N^+(\bA)}
    \left| f^+(mn \cdot (\gamma h, \delta_2 g')) \right| \rd n \rd g_2' \rd h
    \]
is convergent and we have
    \begin{equation}    \label{eq:geo_descent_GL_step_1}
    \int_{[H]} \int_{[G_n']}
    \sum_{ \substack{ \gamma \in P_H(F) \backslash H(F) \\
    \delta_2 \in P_2'(F) \backslash G_n'(F) }}
    k_{f^+, P, \alpha}(\gamma h,\delta_2 g') \eta_{G'}(g') \rd g_2' \rd h =
    k_{\varphi_{f^+}, P, \alpha}(g_1') \eta(g_1').
    \end{equation}
Here we recall that $g' = (g_1', g_2') \in G'(\bA)$, $g_1', g_2' \in G_n'(\bA)$, and $\delta_2 g'$ stands for the element $(g_1', \delta_2 g_2') \in G'(\bA)$. The integration is over $h$ and $g_2'$, and we left with a function on $g'_1$.
\end{lemma}

\begin{proof}
We will prove~\eqref{eq:geo_descent_GL_step_1}. The absolute convergence can
be obtained by the same computation, with only obvious minor modifications.

We first unfold the sum over $\gamma$ and $\delta$ to conclude
that~\eqref{eq:geo_descent_GL_step_1} equals
    \[
    \int_{[H]_{P_H}} \int_{[G_n']_{P_{n}'}}
    k_{f^+, P, \alpha}(h,  g') \eta_{G'}(g') \rd g_2' \rd h.
    \]

We fix good maximal compact subgroups $K_H$ and $K_2'$ of $H(\bA)$ and $G_n'(\bA)$ respectively. By the Iwasawa decomposition the left hand side
of~\eqref{eq:geo_descent_GL_step_1} equals
    \begin{align*}
    \eta(\det g_1')^{n+1} &
    \sum_{m \in M_{\alpha}^+(F)}
    \int_{K_H}
    \int_{K_2'} \int_{[M_{P_H}]} \int_{[M_n']} \int_{N^+(\bA)}
     e^{\langle
    -2\rho_{P_H},H_{P_H}(m_H)\rangle+\langle -2\rho_{P_n'}, H_{P_n'}(m')
    \rangle }
    \\ &f^+ (mn \cdot (m_H k_H, (g_1',m'k')))
     \eta(\det m'k')^{n+1} \rd n \rd m'  \rd m_H  \rd k' \rd k_H.
    \end{align*}
Unfolding the sum over $m \in M_{\alpha}^+(F)$ we conclude that this equals
    \begin{align*}
    &\eta(\det g_1')^{n+1}
    \sum_{m \in M^+_{\alpha}(F)/(M_{P_H}(F)  \times
    M_2'(F))}
    \int_{K_H} \int_{K_2'} \int_{M_{P_H}(\bA)}
    \int_{M_n'(\bA)} \int_{N^+(\bA)}
      \\
    &e^{\langle
    -2\rho_{P_H},H_{P_H}(m_H)\rangle
    +\langle -2\rho_{P_n'},H_{P_n'}(m') \rangle }
    f^+ \left(nm \cdot (m_H k_H, (g_1',m'k'))\right)
    \eta(\det m'k')^{n+1} \rd n \rd m'  \rd m_H \rd k' \rd k_H.
    \end{align*}

We have
    \[
    M^+_{\alpha}(F)/ (M_{P_H}(F)\times
    M_n'(F))=M^+_{n,\alpha}(F)/M_n'(F).
    \]
We also decompose the integral over $N^+(\bA)$ into an integral over ${N_H(\bA)}$ and $N_n^+(\bA)$. An element in $N^+_n(F)$ is considered as an element in $N^+(F)$ via
$(g, x, y) \mapsto ((1, g), x, y)$. Then the above integral equals
    \begin{align*}
    \eta(\det g_1')^{n+1}
    &\sum_{m \in M^+_{n,\alpha}(F)/M_n'(F)}
    \int_{K_H} \int_{K_2'}
    \int_{M_{P_H}(\bA)} \int_{M_n'(\bA)}
    \int_{N_H(\bA)} \int_{N_{n}^+(\bA)}   \\
    & e^{\langle
    -2\rho_{P_H},H_{P_H}(m_H)\rangle+
    \langle -2\rho_{P_n'},H_{P_n'}(m') \rangle }
    f^+ \left(n_H m n_2 \cdot (m_H k_H, (g_1',m'k'))\right)\\
    &\eta(\det m'k')^{n+1} \rd n_2 \rd n_H \rd m' \rd m_H \rd k' \rd k_H.
    \end{align*}
We further break up the integral over $N_{n}^+(\bA)$ into
$N_{n}^+(\bA)/N_n'(\bA)$ and $N_n'(\bA)$ and arrive at
    \begin{equation}    \label{eq:descent_simplification_final_step}
    \begin{aligned}
    \eta(\det g_1')^{n+1}
    &\sum_{m \in M^+_{n,\alpha}(F)/M_n'(F)}
    \int_{K_H} \int_{K_2'}
    \int_{M_{P_H}(\bA)} \int_{M_n'(\bA)}
    \int_{N_H(\bA)} \int_{N_{n}^+(\bA)/N_n'(\bA)} \int_{N_n'(\bA)}   \\
    & e^{\langle
    -2\rho_{P_H},H_{P_H}(m_H)\rangle+
    \langle -2\rho_{P_n'},H_{P_n'}(m') \rangle }
    f^+ \left(n_H m n_2 n_2' \cdot (m_H k_H, (g_1', m'k'))\right)\\
    &\eta(\det m'k')^{n+1}
    \rd n_2' \rd n_2 \rd n_H \rd m' \rd m_H \rd k' \rd k_H.
    \end{aligned}
    \end{equation}
Using the Iwasawa decomposition, we combine the integrals for $m', n'$ and
$k'$ as an integral over $G_n'(\bA)$, and the integrals over $m_H, n_H$ and
$k_H$ as an integral over $H(\bA)$. Note that $\mu(\det mn_2)=1$ for $m \in
M_n'(F)$ and $n \in N_n(\bA)$. We thus conclude
that~\eqref{eq:descent_simplification_final_step} equals
    \[
    \eta(\det g_1') \sum_{m \in M^+_{n,\alpha}(F)/M_n'(F)}
    \int_{N_{n,+}(\bA)/N_n'(\bA)} \varphi_{f^+}
    \left(  g_1'^{-1}\nu(mn_2) g_1' \right)
    \rd n_2.
    \]
Note that there is a slight difference in the definition of $\varphi_{f^+}$
for $n$ even or odd, which eventually simplifies to a uniform expression.
Finally, the map $\nu$ induces a natural bijection between $M_{n,
\alpha}^+(F)/M_n'(F)$ and $M_{n}^+(F) \cap X_{\alpha}(F)$. This concludes
the proof.
\end{proof}

\end{document}
